% Generated draft: requires mathematical and source review.
\documentclass[11pt]{article}
\usepackage[margin=25mm]{geometry}
\usepackage{fontspec}
\setmainfont{Latin Modern Roman}
\usepackage{amsmath,amssymb,mathtools,unicode-math}
\setmathfont{Latin Modern Math}
\usepackage{enumitem,needspace,xurl,hyperref}
\setlength{\parindent}{0pt}
\setlength{\parskip}{5pt}
\setlength{\emergencystretch}{3em}
\newcommand{\N}{\mathbb{N}}
\newcommand{\Z}{\mathbb{Z}}
\newcommand{\Q}{\mathbb{Q}}
\newcommand{\R}{\mathbb{R}}
\newcommand{\C}{\mathbb{C}}
\newcommand{\F}{\mathbb{F}}
\newcommand{\A}{\mathbb{A}}
\newcommand{\PP}{\mathbb{P}}
\newcommand{\E}{\mathbb{E}}
\newcommand{\eps}{\varepsilon}
\newcommand{\dd}{\,\mathrm{d}}
\begin{document}
\section*{UnsolvedMath ID: 700605}
\textbf{Source problem code:} catalogue entry
\subsection{Definitions and mathematical statement}
Let \(K\in\{\R,\C\}\), and let \(\mathrm{SL}_2(K)\) denote the Lie group of determinant-one \(2\times2\) matrices over \(K\). Let \(G_1,G_2\) be Lie groups, each containing a simply connected nilpotent normal Lie subgroup \(B_i\trianglelefteq G_i\), such that
\[
G_i/B_i\cong\mathrm{SL}_2(K)\qquad(i=1,2),
\]
where the isomorphisms are understood as Lie group isomorphisms. Suppose that \(G_1\) and \(G_2\) are embedded as closed subgroups of a topological group \(G\), with their given Lie group topologies, and that
\[
H:=B_1B_2=\{b_1b_2:b_1\in B_1,\ b_2\in B_2\}\subseteq G_1\cap G_2.
\]
In particular, the factors normalize one another, so \(H\) is a subgroup. Assume also that
\[
\text{for every Lie subgroup }L\leq B_1\cap B_2,\qquad
\bigl(gLg^{-1}=L\text{ for every }g\in G\bigr)\Longrightarrow L=\{1\}.
\]
Thus the exclusion concerns every nonidentity Lie subgroup that is normal in the entire ambient group \(G\).

For a group \(H\), define its lower central series by \(\gamma_1(H)=H\) and \(\gamma_{r+1}(H)=[\gamma_r(H),H]\). A bound \(c\) on its nilpotency class means \(\gamma_{c+1}(H)=\{1\}\).

\emph{Question, with uniform boundedness made explicit.} For each \(K\in\{\R,\C\}\), do there exist nonnegative integers \(c_K,d_K\), independent of \(G,G_1,G_2,B_1,B_2\), such that every configuration satisfying these hypotheses obeys
\[
\gamma_{c_K+1}(B_1B_2)=\{1\},\qquad
\dim(B_1B_2)\leq d_K?
\]
Here the dimension refers to the Lie subgroup product. The source does not specify closed versus immersed Lie subgroup conventions, either for this product or for the subgroups \(L\) excluded above. These conventions require confirmation from the original source. The source also does not explicitly state the permitted dependence of the bounds; the formulation above interprets them as uniform over all configurations for each field. Since there are only two fields, bounds for each field would also give common bounds by taking maxima. The source suggests that methods of amalgams from finite group theory might be relevant, without requiring a particular method.
\subsection{Short English statement}
A. L. Chermak asks whether two Lie groups with simply connected nilpotent normal subgroups and quotients \(\mathrm{SL}_2(K)\), for \(K=\R\) or \(\C\), force uniform bounds on the nilpotency class and dimension of \(B_1B_2\) when embedded as specified in a common topological group. The crucial restriction is that no nonidentity Lie subgroup of \(B_1\cap B_2\) is normal in that ambient group. Subgroup conventions and the intended uniformity of the bounds require source confirmation.
\subsection{Sources}
\begin{itemize}
\item Supplied local UnsolvedMath registry record: permanent ID \(700605\), title \emph{Kourovka Notebook Problem 9.72}, exact source problem code \texttt{catalogue entry}, upstream revision \texttt{d371ee917ac9647fbb06cd6fc529deec6f3dda82}.
\item Attributed in the supplied record to A. L. Chermak, proposed in 1984; original problem number 9.72 in \emph{The Kourovka Notebook (Unsolved Problems in Group Theory)}, 9th issue (1984).
\item The supplied background also identifies the 21st issue edition (2026), \texttt{arXiv:1401.0300v48}. The supplied collection description names E. I. Khukhro and V. D. Mazurov as editors. These bibliographic details are supplied attribution and have not been independently verified.
\item The registry labels the problem \texttt{open}. This is a supplied status label, not an independently verified literature assessment. No separate literature-verification date, literature-source links, or page reference is supplied.
\end{itemize}
\end{document}
